\documentclass[11pt]{article}
\usepackage{hyperref}
\usepackage{enumitem}
\usepackage{amsfonts} 
\setlength{\oddsidemargin}{-0.1in}
\setlength{\evensidemargin}{\oddsidemargin}
\setlength{\textwidth}{7.0in}
\setlength{\textheight}{9.3in}
\setlength{\topmargin}{-0.7in}
%\input{macros.tex}
\begin{document}

\fbox{
\vbox{
\begin{flushleft}
Jonathan Llovet \\  % authors' names
COSC 417 \\  %class
2026-04-10\\  % date
\end{flushleft}
\center{\Large{\textbf{Assignment 6}}}
} % end vbox
} % end fbox
\vline

\begin{enumerate}[label=\textbf{Exercise \arabic{*}}]

\item \label{e3}
Show that the language
 \[
B = \{\langle M \rangle \mid \mbox{the Turing machine $M$ accepts at least $3$ strings}\}
\]  is 
 Turing recognizable (Notes: Turing-recognizable sets are also called CE (computably enumerable), and also RE (recursively enumerable);  $\langle M \rangle$ isa standard encoding of the Turing machine $M$). 


\medskip

In your solution you will informally describe an algorithm $S$ that on input $\langle M \rangle$  has the following properties:

(a)	if $M$ accepts at least three strings, then $S$ accepts,

(b)	if $M$ acccepts fewer than three strings, then $S$ rejects or loops.

Of course $S$ will have to run  simulations of $M$ on various inputs,  and the issue is to organize these simulations in the right way to obtain (a) and (b). The main challenge is that $M$ may loop on some inputs, and, therefore, the algorithm $S$ must have a way to avoid being stuck on a simulation of $M$ on an input $x$ on which it loops. Make sure you explain how this challenge is overcome.
\medskip



\textbf{Answer:}. 

\item $A$ and $B$ denote languages, and $\overline{A}$ is the complement of $A$, and $\overline{B}$ is the complement of $B$.  $A \cap B$ is the intersection of $A$ and $B$. Consider the following statements:

(a) $A$ is context-free.

(b) $B$ is Turing-recognizable.

(c) $\overline{A}$ is context-free.

(d)	$\overline{A}$ is decidable.

(e) $A \cap B$ is decidable.

(f) $A \cap B$ is  Turing recognizable.

(g) $\overline{B}$ is Turing recognizable.
\smallskip

Answer the following questions (YES or NO) and provide short justifications:
\smallskip

(1) Does (a) imply (c)?

(2) Does (a) imply (d)?

(3) Does (a) and (b) imply (e)?

(4) Does (a) and (b) imply (f)?

(5) Does (b) imply (g)?
\medskip

\textbf{Answer:}
%\end{enumerate}


\bigskip

\item

Recall that we have shown that the fact ``language $L$  is CE" is equivalent to the existence of a decider  $D $ such that for every $x$, 
\[x \in L \Leftrightarrow \big(\exists~s) ~ D \mbox{ accepts $(x,s)$}\big],
\]
 where  $s$ is a string.
 %\footnote{In Notes 10, we have shown this only for $k=1$, but  that proof can be adapted to work for any constant number $k$, because we can encode a $k$-tuple $(s_1, s_2, \ldots, s_k)$ of numbers into a single number $t$ so that from $t$ one can retrieve each $s_i$ (for example the tuple $(5, 11)$ can be encoded as $t=2^5 3^{11}$).}


Use this to give an alternative proof that the set from~Exercise 1, namely
 \[
B = \{\langle M \rangle \mid \mbox{the Turing machine $M$ accepts at least $3$ strings}\}
\]  is 
 Turing recognizable (or RE). 
For this you need to informally describe the appropriate decider.


\medskip

\textbf{Answer:}


\end{enumerate}




 \end{document}
 